\section{The Bit-Sliced Evaluation Argument}
\label{sec:eval}

This section presents the opening: an interactive (Fiat--Shamir
compiled) argument that certifies the batch of MLE evaluation claims
$a_1,\dots,a_K$ of \Cref{sec:prelim} against the commitment of
\Cref{sec:commitment}. It is a Brakedown/Ligero-style protocol with an
$\Ftwo[X]$-native twist: the evaluation phase produces, per column, a
single degree-$<\Dval$ polynomial over $\K$ whose $\alpha$-projection is
the claim, so the entire argument runs in $\Kpoly$-arithmetic and never
forms an inverse lift.

\subsection{The bit-slice fold}

The technical core is the following observation. The bilinear form
\eqref{eq:claim-bilinear} that defines $a_g$ commutes the eq-tensor
through $\psia$, and $\psia$ is $\Ftwo$-linear; so we may compute the
fold of the \emph{cell bits} first and project last.

\begin{definition}[Bit-slice fold]
\label{def:bitslice}
For a column $w_g$ with cell matrix $M_{w_g}$ and the eq-tensor
$(q_0, q_1)$ of \Cref{sec:eqtensor}, the \emph{bit-slice fold} is the
degree-$<\Dval$ polynomial
\[
  a_g' \;=\; \sum_{b=0}^{\Dval-1} a_g'[b]\,X^b \;\in\; \Kpoly,
  \qquad
  a_g'[b] \;=\; \sum_{i}\sum_{j} q_0[i]\,q_1[j]\,\bigl(M_{w_g}[i,j]\bigr)_b \;\in\; \K,
\]
where $(M_{w_g}[i,j])_b \in \{0,1\}$ is the $b$-th bit of the cell. Its
$b$-th coefficient is the boolean-hypercube fold of the column's $b$-th
bit slice.
\end{definition}

In the codebase $a_g'$ is computed in two passes:
$b_g[i] = \sum_j q_1[j]\cdot M_{w_g}[i,j]$ (a $\K$-scalar times a
$\{0,1\}$-cell, accumulated into $\Kpoly$ by
\code{gf128poly\_accumulate\_cell}, which scatters the scalar $q_1[j]$
into the cell's set-bit coefficient slots), followed by
$a_g' = \sum_i q_0[i]\cdot b_g[i]$ (a $\K$-scalar times a $\Kpoly$,
accumulated by \code{gf128poly\_accumulate\_scaled}).

\begin{lemma}[Bit-slice discharge]
\label{lem:bitslice}
For every column $g$, $\;\psia(a_g') = a_g$.
\end{lemma}

\begin{proof}
By \Cref{rem:psi-extends}, $\psia(a_g') = \sum_b a_g'[b]\,\alpha^b$.
Substituting \Cref{def:bitslice} and exchanging the order of summation,
\[
\begin{aligned}
  \psia(a_g')
  &= \sum_b \Big(\sum_{i,j} q_0[i]q_1[j]\,(M_{w_g}[i,j])_b\Big)\alpha^b
   = \sum_{i,j} q_0[i]q_1[j]\sum_b (M_{w_g}[i,j])_b\,\alpha^b \\
  &= \sum_{i,j} q_0[i]q_1[j]\,\psia(M_{w_g}[i,j]),
\end{aligned}
\]
which is exactly $q_0^\top\,\psia(M_{w_g})\,q_1 = a_g$ by
\eqref{eq:claim-bilinear}.
\end{proof}

\Cref{lem:bitslice} replaces the $\alpha$-lift's central machinery (the
$\Kbits\times\Kbits$ inverse-lift table) with a single inner product:
the claim $a_g$ is read off $a_g'$ as
$a_g = \sum_b a_g'[b]\,\alpha^b$. Crucially, $a_g'$ has $\Dval = 32$
coefficients regardless of the extension degree $\Kbits$, so the
committed claim object is small and independent of the field width.

\subsection{$\gamma$-batching and the proximity fold}

To make the proof constant in the number of columns, the prover draws a
transcript-fresh challenge $\gamma = (\gamma_1,\dots,\gamma_K) \in \K^K$
(one per committed column) and folds all per-column data by it. Define
\begin{align}
  a' &\;:=\; \sum_{g=1}^{K} \gamma_g\,a_g' \;\in\; \Kpoly,
  \label{eq:a-batch}\\
  b'[i] &\;:=\; \sum_{g=1}^{K} \gamma_g\,b_g[i] \;\in\; \Kpoly,
    \qquad i \in [n_{\mathrm{rows}}],
  \label{eq:b-batch}
\end{align}
where $b_g[i] = \sum_j q_1[j]\,M_{w_g}[i,j]$ as above. Separately, after
$\alpha$- and $\gamma$-dependent data is fixed, the prover draws
\emph{proximity coefficients} $c = (c_1,\dots,c_{n_{\mathrm{rows}}}) \in
\K^{n_{\mathrm{rows}}}$ and forms the $\gamma$-batched, $c$-folded
\emph{combined row}
\begin{equation}
  \label{eq:combined-row}
  \mathrm{cr}'[j] \;:=\; \sum_{g=1}^{K} \gamma_g \sum_{i} c_i\,M_{w_g}[i,j]
  \;\in\; \Kpoly,
  \qquad j \in [n_{\mathrm{cols}}].
\end{equation}
The combined row is a random linear combination of the message rows (in
the message space, length $n_{\mathrm{cols}}$); it is the object the
proximity test will encode and check against the committed codewords.

\subsection{The protocol}

\begin{protocol}[Bit-sliced $\gamma$-batched open]
\label{prot:open}
Run on the shared Fiat--Shamir transcript $\mathcal{T}$ that already
carries the commitment and the upstream challenges $(\alpha, r^*)$ and
claims $(a_g)$. Let $t = \code{num\_column\_openings}$ and
$N = n_{\mathrm{cols}}\cdot\mathrm{REP}$.

\smallskip
\noindent\textbf{Prover.}
\begin{enumerate}[leftmargin=2em]
  \item Compute the eq-tensor $(q_0, q_1)$ from $r^*$ (kept in $\K$).
  \item Draw $\gamma \leftarrow \mathcal{T}$ ($K$ challenges).
  \item For each column $g$: compute $b_g$ and $a_g'$
    (\Cref{def:bitslice}); accumulate $a' \mathrel{+}= \gamma_g a_g'$ and
    $b'[i]\mathrel{+}= \gamma_g b_g[i]$.
  \item Absorb $b'$ then $a'$ into $\mathcal{T}$.
  \item Draw the proximity coefficients $c \leftarrow \mathcal{T}$
    ($n_{\mathrm{rows}}$ challenges).
  \item Compute the combined row $\mathrm{cr}'$ \eqref{eq:combined-row};
    absorb it into $\mathcal{T}$.
  \item For $k = 1,\dots,t$: draw a column index $j_k \leftarrow
    \mathcal{T}$ (uniform in $\{0,\dots,N-1\}$ by masking a squeezed
    word); emit the opened codeword column $\bigl(\mat{C}_g[i,j_k]\bigr)_{g,i}$
    and its Merkle authentication path $\pi_k$.
  \end{enumerate}
The proof is $(a',\, b',\, \mathrm{cr}',\, \{(j_k,\,\text{cells}_k,\,\pi_k)\}_{k=1}^t)$.

\smallskip
\noindent\textbf{Verifier.}
Re-derive $(q_0,q_1)$, $\gamma$, $c$, and the indices $j_k$ from
$\mathcal{T}$ in the prover's order (absorbing $b',a'$ before $c$, and
$\mathrm{cr}'$ before the $j_k$), then run \Cref{cons:checks}.
\end{protocol}

\begin{construction}[The four checks]
\label{cons:checks}
The verifier accepts iff all four hold; the first failure rejects.
\begin{enumerate}[leftmargin=3.4em, label=\emph{Check \arabic*.}, ref=Check~\arabic*]
  \item\label{chk:eval} \emph{(Evaluation consistency.)}
    \[
      \sum_{i} q_0[i]\cdot b'[i] \;\stackrel{?}{=}\; a'
      \qquad\text{in } \Kpoly.
    \]
  \item\label{chk:discharge} \emph{($\alpha$-projection discharge.)}
    \[
      \psia(a') \;=\; \sum_{b} a'_b\,\alpha^b
      \;\stackrel{?}{=}\; \sum_{g=1}^{K} \gamma_g\,a_g
      \qquad\text{in } \K.
    \]
  \item\label{chk:coherence} \emph{(Coherence.)}
    \[
      \sum_{j} q_1[j]\cdot \mathrm{cr}'[j]
      \;\stackrel{?}{=}\;
      \sum_{i} c_i\cdot b'[i]
      \qquad\text{in } \Kpoly.
    \]
  \item\label{chk:encoding} \emph{(Encoding consistency at each opened
    column.)} For every opened index $j_k$:
    \begin{itemize}[leftmargin=2em, label=---]
      \item \emph{Merkle:} $\pi_k$ verifies the cells
        $\bigl(\mat{C}_g[i,j_k]\bigr)_{g,i}$ against $\mathrm{root}$ at leaf
        $j_k$.
      \item \emph{Encoding identity:}
        \[
          \enc(\mathrm{cr}')[j_k]
          \;\stackrel{?}{=}\;
          \sum_{i} c_i \sum_{g=1}^{K} \gamma_g\,\mat{C}_g[i,j_k]
          \qquad\text{in } \Kpoly,
        \]
        where $\enc$ is the $\Kpoly$-extended RAA encoder of
        \Cref{lem:linearity} and each $\mat{C}_g[i,j_k]\in\cellring$ is
        read as a $\{0,1\}$-cell scattered by the $\K$-scalar $\gamma_g$.
    \end{itemize}
\end{enumerate}
\end{construction}

\begin{remark}[Reading the checks]
\ref{chk:eval} ties $a'$ to the row-folds $b'$; \ref{chk:discharge} ties
$a'$ to the sumcheck claims; \ref{chk:coherence} ties the proximity fold
$\mathrm{cr}'$ to the same $b'$; \ref{chk:encoding} ties $\mathrm{cr}'$ to
the committed codewords. The chain
$a_g \leftrightarrow a' \leftrightarrow b' \leftrightarrow
\mathrm{cr}' \leftrightarrow \{\mat{C}_g\}$ closes the loop from the
claim to the commitment. Only \ref{chk:discharge} leaves $\Kpoly$: it is
the single $\K$-valued projection of the whole argument, and the one
place the field structure of $\K$ (not just its $\Ftwo$-vector-space
structure) is used.
\end{remark}

\begin{remark}[Well-typedness]
Every quantity in \ref{chk:eval}, \ref{chk:coherence}, and
\ref{chk:encoding} is an element of $\Kpoly$: products of a
$\K$-scalar with a $\Kpoly$ stay degree-$<\Dval$, and sums are
coefficientwise; equality is exact, with no reduction.
\ref{chk:discharge} is an equality in $\K$. This is in deliberate contrast to the $\alpha$-lift
predecessor, where the analogous quantities were unreduced $\Ftwo[X]$
polynomials of growing degree ($\approx \Dval + 3\Kbits$ bits for the
batched $a'$); the bit-sliced form keeps a fixed degree bound $\Dval$
because the coefficients absorb the field arithmetic.
\end{remark}

\subsection{Virtual and public columns}
\label{sec:virtuals}

The commitment of \Cref{sec:commitment} stores only the \emph{primary
witness} columns. The full trace, however, contains two further kinds of
column that the argument must still account for without enlarging the
Merkle tree. Both are handled by the verifier folding their contributions
into the existing checks (\code{verify\_f2\_open\_with\_virtuals}); the
batching vector $\gamma$ and the per-column claims $(a_g)$ range over the
\emph{whole} witness trace, so $K$ below counts primary and virtual
columns alike, while the committed codewords $\{\mat{C}_g\}$ exist only
for the primary ones.

\paragraph{Bit-operation virtuals.}
A trace column $w_h$ may be declared as a fixed bitwise map of a committed
\emph{source} column $w_s$, applied cell by cell:
$w_h[\cdot] = T\!\bigl(w_s[\cdot]\bigr)$, where $T$ is a cyclic rotation
$\mathrm{ROTL}_c$ or a logical right shift $\mathrm{SHR}_c$ of the
$\Dval$-bit cell (codebase \code{BitOp::Rot}\,/\,\code{BitOp::ShiftR},
realised by \code{apply\_bit\_op\_u32}). Such a column is not committed.
This is sound because every such $T$ is an $\Ftwo$-linear map on the
cell's bit-vector --- a rotation is a coordinate permutation, a shift is a
coordinate truncation, both $0/1$ matrices on $\Ftwo^{\Dval}$ --- and
therefore, by exactly the argument of \Cref{lem:linearity}, commutes with
the RAA encoder applied cell-wise:
\begin{equation}
  \label{eq:virtual-commute}
  \enc\!\bigl(T(w_s[i,\cdot])\bigr)
  \;=\; T\!\bigl(\enc(w_s[i,\cdot])\bigr)
  \;=\; T\!\bigl(\mat{C}_s[i,\cdot]\bigr).
\end{equation}
(Repetition and the permutations are cell-positional and trivially
commute with the per-cell $T$; the prefix-XOR accumulator commutes because
$T(a)+T(b)=T(a+b)$ in $\Ftwo$.) Thus the virtual column's codeword is
\emph{determined} by --- and recomputable from --- the source's committed
codeword: it needs no commitment of its own.

The virtual column nonetheless participates in the argument like any
other, carrying its own batching scalar $\gamma_h$ and its own MLE claim
$a_h$ (delivered by the sumcheck for $w_h$ as a genuine trace column):
\begin{itemize}[leftmargin=2em]
  \item In \ref{chk:discharge} it contributes the term $\gamma_h\,a_h$ to
    the right-hand side, indistinguishably from a committed column.
  \item In the prover's folds $a'$, $b'$, and $\mathrm{cr}'$ the column is
    materialised, so its contribution is genuinely present in the
    transmitted $\mathrm{cr}'$.
  \item In \ref{chk:encoding} the verifier, assembling the per-opened-column
    fold $\sum_i c_i\sum_g \gamma_g\,\mat{C}_g[i,j_k]$, has \emph{no}
    committed cell $\mat{C}_h[i,j_k]$ to read. It instead reconstructs that
    cell as $T\!\bigl(\mat{C}_s[i,j_k]\bigr)$ from the \emph{source's}
    opened cell --- which the Merkle path \emph{does} certify --- and
    scatters it by $\gamma_h$. Identity \eqref{eq:virtual-commute} is
    precisely what makes this reconstruction equal the prover's
    materialised contribution to $\enc(\mathrm{cr}')[j_k]$, so the encoding
    identity still closes.
\end{itemize}
Two restrictions keep this self-consistent and sound. The derivation is
\emph{depth one}: a source column must itself be primary, never another
virtual (enforced by \code{primary\_col\_indices\_from\_virtuals}), so a
single opened source cell suffices and no chain of reconstructions arises.
And $T$ must be $\Ftwo$-linear: nonlinear maps such as $\mathrm{AND}$ or
$\mathrm{Maj}$ do \emph{not} commute with the encoder and cannot be
virtualised this way; the host protocol discharges those instead through
the second functional read-off $\psiz$ developed in
\Cref{sec:multifunctional}.

\paragraph{Public columns.}
Public inputs are not committed at all. The verifier evaluates their MLEs
at $r^*$ directly from the public data and omits them from the opened
batch; they enter neither $\gamma$ nor the codeword fold. This mirrors the
witness-only commitment of \Cref{rem:storage}.

\medskip
Neither refinement alters the algebra of \Cref{cons:checks}: they only
redefine which terms enter the sums --- a virtual column substitutes
$T(\mat{C}_s)$ for a committed $\mat{C}_h$, and a public column drops out.
